\chapter{How to read}

As with any workbook, i.e.~book that's there to guide actual work not to teach a
big subject, there are numerous ways how to read the text. 
You can read the manuscript chapter by chapter, but this is unlikely to be the
most fun experience.
Below, we collect some ideas how to substructure the reading:


\section*{Performance Analysis Course}

We have made reasonable good experience with the following order of activities:

\begin{enumerate}
  \item Preparation
  \begin{itemize}
    \item Ask participants to prepare a benchmark following the recommendations
     in Section~\ref{section:preparation:benchmarks}.
  \end{itemize}
  \item First session
  \begin{itemize}
    \item Use the first hour or so to briefly run through the compiler settings
  in Section~\ref{section:preparation:compiler-settings} to ensure that people
  do not gather biased/invalid data. In this context, it might also be
  reasonable to ask them to disable I/O (Section~\ref{section:preparation:io}).
    \item Start straight with the high-level assessment from
      Chapter~\ref{chapter:high-level}. Throughout this exercise, participants
      will learn about the machinery used on-the-fly
      (cmp.~Section~\ref{section:preparation:machine}).
  \end{itemize}
  \item Second session
  \begin{itemize}
    \item Run with participants through some basic performance analysis
    terminology from Section~\ref{section:terminology:hpc-terms}, 
    \item and then introduce the concept of localisation
    (Section~\ref{section:tools:localisation}).
    \item Discuss a range of profilers from Chapter~\ref{chapter:tools} (we
    always do at least gprof, Map and VTune) and let people expore their code
    with these tools.
  \end{itemize}
  \item Core sessions
  \begin{itemize}
    \item Run through the individual chapters starting from
    Chapter~\ref{chapter:core} one by one.
  \end{itemize}
  \item Wrap-up
  \begin{itemize}
    \item Let people read through the case studies in
    Part~\ref{part:case-studies}.
    \item Ask the audience to present their lessons learned.
  \end{itemize}
\end{enumerate}



\section*{First-time assessment}

We recommend to run through the preparatory remarks in
Chapter~\ref{chapter:preparation} first and to complete all the checkboxes (dos)
there.
Once completed, it is time to run through the high-level assessment remarks in
Chapter~\ref{chapter:high-level}.
With most codes, it is highly unlikely that a code will only show one kind of
flaw in the high-level analysis. 

We propose to run through the detailed analysis step by step afterwards.
Per step, it makes sense to consult the tools summary
(Chapter~\ref{chapter:tools}) and notably to try out different things.
Where appropriate, the reader might prefer to read through one of the case
studies in Part~\ref{part:case-studies}, just to get some inspiration what tools
and approaches to try out and how to report outcomes.



\section*{Assessment}

If you have already completed numerous assessments, the the manuscript is
mainly a reference document.
You'd start with the high-level assessment in Chapter~\ref{chapter:high-level},
before you dig into the various subchapters depending on the initial assessment
outcome.
The case studies are not of particular relevance (anymore), as you have already
completed you own case studies.


\section*{Feedback}

The quintessential objective of any assessment is to give feedback to developers.
Once they make changes to their source code, it is important to re-start the 
assessment from scratch.
There is always this temptation to return to the particular characteristics of interest that have led to a report, to see how the code changes by the development team have altered these characteristics, and then to continue from there.


It is important to reassess the ``flaw'' metrics.
However, this should only be a validation that something has changed.
After that, we have to start from scratch, as performance optimisation is a complicated process where one change might introduce flaws, improvement or character changes in a completely different part of the code or affect a completely different flavour of the performance.
If we alter a particular code part that makes the scaling deteriorate for example, this might have a severe knock on effect on the core performance.
If we improve the core performance, the GPU offloading might all at a sudden run into data transfer issues.
There are numerous examples for these side-effects.


\paragraph*{How it works}

\begin{itemize}[noitemsep]
  \item[$\boxplus$] Conduct a performance assessment which both covers all aspects of the code (cmp.~Chapter~\ref{chapter:high-level}) plus goes in-depth all the way through for all flaws identified.
  \item[$\boxplus$] Hand this completed assessment back to the developers to address the flaws.
  \item[$\boxplus$] Restart the assessment from scratch top-down, i.e.~starting with the high-level assessment covering all code runtime nuances.
\end{itemize}


\paragraph*{How it does not work}

\begin{itemize}[noitemsep]
  \item[$\boxminus$] Mix code optimisation and performance assessment workflows.
  \item[$\boxminus$] Restart the assessment from the place where you stopped when you reported a flaw and continue digging through the assessment decision metrics without redoing a high-level, overarching assessment.
\end{itemize}
